\documentclass[12pt]{article}
\usepackage[brazil]{babel}
\usepackage[utf8]{inputenc}
\usepackage[T1]{fontenc}
\usepackage{amsmath, amssymb}
\usepackage{geometry}
\usepackage{graphicx}
\usepackage{xcolor}
\usepackage{enumitem}

\definecolor{maincolor}{RGB}{0,100,50}
\usepackage{fancyhdr}
\definecolor{forestgreen}{RGB}{0,100,50}
\usepackage[most]{tcolorbox}
\newtcolorbox{block}[1]{%
  colback=green!5,
  colframe=maincolor,
  coltitle=white,
  title=#1,
  fonttitle=\bfseries,
  boxrule=1.5pt,
  arc=3mm,
  left=6pt,
  right=6pt,
  top=6pt,
  bottom=6pt
}

\geometry{margin=2.5cm}
\pagestyle{fancy}
\fancyhf{}
 
\renewcommand{\headrulewidth}{2pt}
\renewcommand{\footrulewidth}{2pt}

\renewcommand{\headrule}{\hbox to\headwidth{%
  \color{forestgreen}\leaders\hrule height \headrulewidth\hfill}}

\renewcommand{\footrule}{\hbox to\headwidth{%
  \color{forestgreen}\leaders\hrule height \footrulewidth\hfill}}

\setlength{\headheight}{32pt}

\fancyhead[L]{%
  \includegraphics[height=1.2cm]{Universidade_de_Rio_Verde.jpg}
}

\fancyhead[R]{%
  \raggedleft
  \textcolor{forestgreen}{%
    \small
    \textbf{Lista de Exercícios — Relações}\\
    Me.~Sergio Souza Novak\\
    \textit{Lógica para Engenharia de Software}
  }
}

\fancyfoot[C]{%
  \textcolor{forestgreen}{\thepage}
}

% Caixa para resposta discursiva
\newcommand{\resposta}[1]{
\noindent
\framebox[\textwidth]{
\begin{minipage}{0.96\textwidth}
\vspace{#1}
\end{minipage}
}
}

\begin{document}
\vspace*{0.2cm}

\section*{Instruções}
\begin{block}{Definições importantes}
Uma \textbf{relação binária} $R$ em um conjunto $A$ é um subconjunto de $A \times A$.

\textbf{Reflexiva:} $\forall a \in A,\; (a,a) \in R$.

\textbf{Simétrica:} $\forall a, b \in A,\; (a,b) \in R \Rightarrow (b,a) \in R$.

\textbf{Transitiva:} $\forall a, b, c \in A,\; ((a,b) \in R \land (b,c) \in R) \Rightarrow (a,c) \in R$.

\textbf{Relação de equivalência:} é uma relação reflexiva, simétrica e transitiva.
\end{block}
Justifique rigorosamente todas as respostas discursivas.

\section*{Exercícios Discursivos}

\begin{enumerate}

% ============================================================
% QUESTÃO 1 — Discursiva
% ============================================================
\item Seja $A = \{1, 2, 3\}$ e $R = \{(1,1), (2,2), (3,3), (1,2), (2,1)\}$. Verifique se $R$ é reflexiva, simétrica e transitiva. Conclua se $R$ é uma relação de equivalência.

\resposta{6cm}

\vspace{0.8cm}

% ============================================================
% QUESTÃO 2 — Discursiva
% ============================================================
\item Considere a relação $R$ em $\mathbb{Z}$ definida por: $aRb \Leftrightarrow a - b$ é múltiplo de 4. Prove que $R$ é uma relação de equivalência.

\resposta{8cm}

\vspace{0.8cm}

% ============================================================
% QUESTÃO 3 — Discursiva
% ============================================================
\item Seja $A = \{1, 2, 3, 4\}$ e $R = \{(a, b) \in A \times A : a \leq b\}$. Determine se $R$ é reflexiva, simétrica e transitiva. Justifique cada propriedade.

\resposta{7cm}

\vspace{0.8cm}

% ============================================================
% QUESTÃO 4 — Discursiva
% ============================================================
\item Considere a relação $R$ em $\mathbb{Z}$ definida por: $aRb \Leftrightarrow |a| = |b|$. Verifique, com demonstração, se $R$ é reflexiva, simétrica e transitiva.

\resposta{7cm}

\vspace{0.8cm}

% ============================================================
% QUESTÃO 5 — Discursiva
% ============================================================
\item Seja $A = \{1, 2, 3, 4, 5\}$ e a relação $R = \{(1,1),(2,2),(3,3),(4,4),(5,5),(1,3),(3,1),(2,4),(4,2)\}$. Verifique se $R$ é uma relação de equivalência. Caso não seja, indique quais propriedades falham e dê contraexemplos.

\resposta{7cm}

\newpage

\section*{Exercícios de Múltipla Escolha}

% ============================================================
% QUESTÃO 6 — Múltipla escolha
% ============================================================
\item Considere as relações binárias em $\mathbb{Z}$:
\begin{itemize}
    \item $R_1 = \{(a, b) : a + b \text{ é par}\}$;
    \item $R_2 = \{(a, b) : a < b\}$;
    \item $R_3 = \{(a, b) : a \cdot b \geq 0\}$;
    \item $R_4 = \{(a, b) : a - b \text{ é múltiplo de 3}\}$.
\end{itemize}

São relações de equivalência apenas:
\begin{enumerate}[label=\Alph*)]
    \item $R_1$ e $R_2$.
    \item $R_1$ e $R_4$.
    \item $R_2$ e $R_3$.
    \item $R_1$, $R_3$ e $R_4$.
    \item $R_3$ e $R_4$.
\end{enumerate}

\vspace{0.8cm}

% ============================================================
% QUESTÃO 7 — Múltipla escolha
% ============================================================
\item Seja $A = \{1,2,3,4\}$ e $R = \{(1,1),(2,2),(3,3),(4,4),(1,2),(2,1),(3,4),(4,3),(1,3),(3,1)\}$.

Assinale a alternativa correta sobre $R$.
\begin{enumerate}[label=\Alph*)]
    \item $R$ é reflexiva e simétrica, mas não é transitiva.
    \item $R$ é uma relação de equivalência.
    \item $R$ é reflexiva e transitiva, mas não é simétrica.
    \item $R$ é apenas reflexiva.
    \item $R$ não é reflexiva.
\end{enumerate}

\vspace{0.8cm}

% ============================================================
% QUESTÃO 8 — Múltipla escolha
% ============================================================
\item Considere as relações em $\mathbb{Z}$:
\begin{itemize}
    \item $R_1 = \{(a,b) : a \mid b\}$ ($a$ divide $b$);
    \item $R_2 = \{(a,b) : a - b \text{ é múltiplo de 2}\}$;
    \item $R_3 = \{(a,b) : a \neq b\}$.
\end{itemize}

Assinale a alternativa correta.
\begin{enumerate}[label=\Alph*)]
    \item $R_1$ é reflexiva e transitiva, mas não simétrica.
    \item $R_2$ e $R_3$ são relações de equivalência.
    \item $R_3$ é simétrica e transitiva.
    \item $R_1$ é uma relação de equivalência.
    \item Nenhuma das três é reflexiva.
\end{enumerate}

\vspace{0.8cm}

% ============================================================
% QUESTÃO 9 — Múltipla escolha
% ============================================================
\item Seja $A = \{a, b, c\}$ e considere a relação $R = \{(a,a),(b,b),(c,c),(a,b),(b,a),(b,c),(c,b)\}$.

Assinale a alternativa correta.
\begin{enumerate}[label=\Alph*)]
    \item $R$ é uma relação de equivalência.
    \item $R$ é reflexiva e simétrica, mas não é transitiva.
    \item $R$ é reflexiva, mas não é simétrica nem transitiva.
    \item $R$ é simétrica e transitiva, mas não é reflexiva.
    \item $R$ é apenas simétrica.
\end{enumerate}

\vspace{0.8cm}

% ============================================================
% QUESTÃO 10 — Múltipla escolha
% ============================================================
\item Sobre relações binárias em um conjunto $A$, analise as afirmações:

\begin{enumerate}[label=\Roman*.]
    \item Toda relação reflexiva e simétrica é necessariamente transitiva.
    \item Se $R$ é simétrica e transitiva, e para todo $a \in A$ existe $b$ tal que $aRb$, então $R$ é reflexiva.
    \item A relação vazia em um conjunto não vazio é simétrica e transitiva, mas não é reflexiva.
    \item Se $R$ é transitiva e $aRb$, então $aRa$.
\end{enumerate}

Estão corretas apenas:
\begin{enumerate}[label=\Alph*)]
    \item I e II.
    \item II e III.
    \item I, II e IV.
    \item III e IV.
    \item I e III.
\end{enumerate}

\end{enumerate}

\end{document}
